
\section{Robot Design}

With respect to robot design, there are many factors to consider, such as the robot morphology, type of actuators, sensor suite, and compute capabilities.

Morphology refers to the physical structure of the robot as defined by the number of limbs, joints, and their kinematic alignments. For multi-legged robots, two main morphologies are commonly used: parasagittal and sprawling (see figures \ref{fig:cheetah3} and \ref{fig:thex}). Each of these morphologies can have a varying number of legs, with quadrupeds and hexapods being the most common. Both morphologies have their own advantages and disadvantages. Most prominently, parasagittal robots are more energy efficient, with legs straight below the hip, allowing the robot to stand without exerting torque on the joints. However, they usually have a higher center of mass and a smaller support polygon, making them more susceptible to tipping over during locomotion. On the other hand, sprawling robots, with their legs extending in the lateral direction from the hip, have a much lower center of mass and a wider support polygon, making them more stable, but less energy efficient as the joints must exert torque to hold the body up against gravity even when the robot is standing still.

% \begin{figure}[h!]
%     \begin{minipage}[t]{0.48\textwidth}
%         \centering
%         \includegraphics[width=\textwidth]{figures/cheetah3.png}
%         \caption{MIT's Cheetah 3, a parasagittal quadruped with 3DoF legs \cite{bledtMITCheetah32018}}
%         \label{fig:cheetah3}
%     \end{minipage}
%     \hfill
%     \begin{minipage}[t]{0.48\textwidth}
%         \centering
%         \includegraphics[width=\textwidth]{figures/THex.jpeg}
%         \caption{Lynxmotion's T-Hex, a sprawling hexapod with 3DoF legs \cite{THex3DoFHexapod}}
%         \label{fig:thex}
%     \end{minipage}
%     % \caption{Overall caption for both images (optional).}
% \end{figure}

Some examples of state-of-the-art parasagittal robots include the MIT Cheetah 3 \cite{bledtMITCheetah32018}, Boston Dynamics' Spot \cite{Spot}, ETH Zurich's SpaceBok \cite{armSpaceBokDynamicLegged2019}, ANYbotics' ANYmal \cite{hutterANYmalHighlyMobile2016}, IIT's HyQ \cite{seminiHyQDynamicLocomotion}, Unitree B2 \cite{UnitreeB2}, Unitree Go 2 \cite{UnitreeGo2}. All of these robots are quadrupeds having 3 DOF per leg, except for SpaceBok with 2 DOF per leg. On the other hand, examples of sprawling robots include CSIRO's Weaver \cite{bjelonicWeaverHexapodRobot2018}, CTU Prague's HAntR \cite{cizekDesignConstructionRoughTerrain2021}, and Lynxmotion's T-Hex \cite{THex3DoFHexapod}, which have 5-DOF, 4-DOF, and 3-DOF legs respectively. While parasagittal quadrupeds and sprawling hexapods are the most common morphologies, other morphologies such as parasagittal hexapods, for example DOBOT's Hexplorer \cite{DOBOTHEXAPODRobot} and RHex \cite{PDFRHexSimple}, and sprawling quadrupeds, for example the TITAN-XIII \cite{PDFTITANXIIISprawlingtype}, have also been developed.

Beyond morphology, the actuation scheme fundamentally dictates a robot's dynamic capabilities and its ability to negotiate rough terrain. High-performance platforms rely on torque-controlled actuators that allow the robot to act as a virtual spring, absorbing unexpected impacts from uneven ground. For example, the MIT Cheetah 3 \cite{bledtMITCheetah32018} utilizes custom quasi-direct drive (QDD) actuators, which combine high-torque brushless motors with low-ratio gearboxes to achieve high mechanical transparency and backdrivability. ANYmal \cite{hutterANYmalHighlyMobile2016} and StarlETH \cite{hutterSTARLETHCOMPLIANTQUADRUPEDAL2012} employ Series Elastic Actuators (SEAs), placing a physical spring between the motor and the joint to absorb shocks and store energy, while the HyQ platforms \cite{seminiHyQDynamicLocomotion} rely on hydraulic cylinders for extreme force density and impact resistance. Similarly, the industrial-grade Unitree B2 \cite{UnitreeB2} utilizes high-torque-density motors capable of supporting up to 120 kg payloads over severe terrain. To bridge the gap to more accessible hardware, open-source research platforms like PADWQ \cite{kimPADWQDesignControl2021} achieve similar dynamic transparency by leveraging custom low-gear-ratio transmissions and 3D-printed components. Conversely, low-cost and educational platforms, such as the Stanford Pupper \cite{YerbertDingoQuadrupedBase}, HAntR \cite{cizekDesignConstructionRoughTerrain2021}, and our T-Hex kit \cite{THex3DoFHexapod}, utilize position-controlled hobby servos. While affordable and lightweight, these rigid actuators lack force feedback (and often, any kind of feedback) and cannot naturally absorb shocks, requiring specialized algorithms to detect and react to contacts.

To traverse rough terrain and assess spatial constraints, the robot needs a suite of sensors, which is generally divided into proprioceptive and exteroceptive sensors. Proprioceptive sensors measure the robot's internal state; almost all modern legged robots feature high-resolution joint encoders to measure joint positions, torque sensors to measure joint efforts, and a central Inertial Measurement Unit (IMU) placed near the Center of Mass (CoM) to estimate body orientation, position, acceleration, and linear and angular velocity. Contact sensors, detecting when the feet make and break contact with the ground, are also common. Relying heavily on this high-frequency proprioceptive feedback, platforms such as the Ghost Robotics Vision 60 \cite{Vision60Ghost} emphasize blind locomotion to traverse tall grass or deep water where cameras might fail. Similarly, the MIT Cheetah 3 \cite{bledtMITCheetah32018} utilizes proprioception and robust sensor fusion \cite{bledtContactModelFusion2018} to run and gallop at high speeds, bounding up stairs and traversing over unstructured terrain completely blind. Highly specialized robots like SpaceBok \cite{armSpaceBokDynamicLegged2019} also rely heavily on internal state estimation, integrating precise IMUs and leg kinematics to achieve in-flight attitude control in low-gravity environments.

However, exteroceptive sensors, which measure the outside world, are critical for posture adaptation in confined spaces. State-of-the-art robots like Spot \cite{Spot, boumanAutonomousSpotLongRange2020} and ANYmal \cite{gehringANYmalFieldSolving2021} feature an array of stereo depth cameras placed around the perimeter of the chassis to provide 360-degree close-range obstacle detection, often coupled with a 3D LiDAR at the top of the chassis for long-range spatial mapping. Other platforms achieve environmental awareness through more directional sensing. To achieve similar perceptive capabilities on a budget, platforms like the Unitree Go2 \cite{UnitreeGo2} and PADWQ \cite{kimPADWQDesignControl2021} mount a single forward-facing depth camera (such as an Intel RealSense) or a solid-state LiDAR whereas Weaver \cite{bjelonicWeaverHexapodRobot2018} utilizes stereo cameras to perceive the terrain. 

Finally, running advanced locomotion algorithms (the kinds of which we will discuss in the next section) that take in sensor data and output joint commands requires significant onboard compute capabilities, typically arranged in a hierarchical architecture. The low-level control loop, which is mainly responsible for the actual locomotion itself, that is estimating states, calculating joint commands to move and maintain balance, must run at high frequencies (often between 50 Hz and 1000 Hz) and is usually handled by a dedicated real-time microcomputer. The MIT Cheetah 3 \cite{bledtMITCheetah32018} runs its purely proprioceptive locomotion controller at 1000Hz using a dedicated onboard x86 PC running a real-time operating system. For systems with perception modules, high-level planning tasks, such as processing dense point clouds to localize the robot, assess spatial constraints, build a map, or plan a trajectory run at much lower frequencies (5 Hz to 50 Hz), mainly due to computational expenses, but also due to assumptions regarding how fast the perception data realistically changes. ANYmal \cite{hutterANYmalHighlyMobile2016} employs multiple onboard Intel Core i7 processors, dedicating one exclusively to real-time locomotion control and another entirely to perceptive navigation and mapping \cite{hoellerANYmalParkourLearning2023}. PADWQ \cite{kimPADWQDesignControl2021} explicitly incorporates an onboard GPU to handle its perceptive locomotion.

Considering its sprawling morphology, position-controlled servos with less than 1 Nm peak torque, no integrated feedback, and a low-speed servo controller, the T-Hex, especially with two legs removed, is greatly limited by its hardware as compared to other robots in the similar weight class, requiring us to augment it with the bare-minimum sensors for robust locomotion, and work around the limitations of its actuators and servo controller. Algorithms that require minimal state feedback, such as joint state and body orientation, and can run at low frequencies are most suitable for the T-Hex. This motivates an extensive literature review of locomotion approaches to find the most suitable ones for our hardware.

\section{Locomotion and Perception}

The locomotion of legged robots in various challenging scenarios has been an active area of research for the past few decades. Locomotion controllers for legged robots can be divided into two-main categories: model-free controllers and model-based controllers. Here, the model being referred to is the dynamics model of the robot, which is typically used to relate the forces and torques on the robot's legs to the motion of the robot body\cite{bledtMITCheetah32018, dicarloDynamicLocomotionMIT2018}. 

One of the simplest model-free control techniques to make a legged robot walk is to use an open-loop\footnote{As opposed to closed-loop, which means to incorporate state feedback.} kinematic gait that has each leg follow a pre-computed foot trajectory\cite{sunHexapodRobotKinematics2017}. A controller like this is easy to implement thanks to well-established approaches for deriving legs' kinematic equations, but is expected to perform poorly in anything but flat terrain. A legged robot is a highly dynamic system, and such a controller ignores the dynamics entirely. Furthermore, not incorporating feedback means the controller is blind to when a leg doesn't reach the desired position. If the leg hits an obstacle in the middle of its swing phase blocking its movement, the leg will continue to push against the obstacle, instead of adjusting its trajectory, like by increasing the amplitude of the swing, to overcome it. These factors, combined with the reliance on good gait parameters such as swing amplitude, stride length, and the lateral distance from the hip\footnote{This is only applicable for robots with a sprawling morphology, like T-Hex. For parasagittal robots, such as Spot or ANYmal, the lateral distance from the hip can be 0.}, makes such a controller unsuitable for unstructured terrain. 

Central Pattern Generators (CPGs), inspired by biological neural networks in animal nervous systems, are used in more sophisticated model-free controllers to generate kinematic references for the robot's joints to follow. The implementation of CPGs varies across literature \cite{owakiQuadrupedRobotExhibiting2017, barasuolReactiveControllerFramework2013, fukuokaSimpleRuleQuadrupedal2015, AdaptiveDynamicWalking, chenLocomotionControlSensordriven2016, parkStableModifiableWalking2013, zhangTrotGaitDesign2014}, but at a high-level, a CPG is attached to each joint in the robot and generates oscillatory signals that act as reference joint trajectories. CPGs in each joint can communicate with each other implicitly through mathematical coupling, which allows for designing specific gaits that rely on phase differences between groups of legs. Performance of CPGs also relies on the choice of parameters, and learning-based methods or derivative-free optimization techniques like evolutionary algorithms are often used to find the best set of parameters by simulating the robot in virtual terrain \cite{liParameterTuningCPGs2015, renMultipleChaoticCentral2015, endoLearningCPGbasedBiped2008}. When state feedback is incorporated into such controllers, they can be surprisingly robust even in non-flat terrain involving slopes, stairs, and unstructured elements such as pebbles and stones \cite{AdaptiveDynamicWalking, barasuolReactiveControllerFramework2013, chenLocomotionControlSensordriven2016}.

Recently, Deep Reinforcement Learning (DRL) has been under the spotlight for learning robust locomotion controllers for legged robots. Proximal Policy Optimization (PPO), a policy-gradient method developed at OpenAI \cite{schulmanProximalPolicyOptimization2017}, is used when training DRL-based controllers due to its impressive performance particularly for tasks related to robot locomotion and manipulation. In PPO, two separate neural networks are trained simultaneously. The value network is responsible for learning how good each state is, whereas the policy network learns the best actions to take at each state - only the latter is actually used after training. To train the two networks, many instances of robots are simulated in virtual environments \cite{rudinLearningWalkMinutes2022}. NVIDIA's Isaac Lab \cite{nvidiaIsaacLabGPUAccelerated2025} provides pre-developed plug-and-play packages for training the locomotion of any legged robot in a variety of environments. To minimize the sim-to-real gap, domain randomization is used to simulate unmodeled factors and wear and tear of hardware, which makes the controller robust to external disturbances. DRL-based locomotion controllers are successful in a wide variety of terrains and demonstrate impressive dynamic movement abilities \cite{tahaDevelopingHexapodLocomotion2025, mikiLearningRobustPerceptive2022, leeLearningQuadrupedalLocomotion2020, choiLearningQuadrupedalLocomotion2023, margolisRapidLocomotionReinforcement2022, tanSimtoRealLearningAgile2018, rudinCatlikeJumpingLanding2022, rudinAdvancedSkillsLearning2022, nahrendraDreamWaQLearningRobust2023, kimDeepReinforcementLearning2026}.

Model-based approaches require some dynamic model of the robot. A perfectly accurate dynamic model is usually too complex to be directly used; however, under specific assumptions, a simplified dynamic model can be used to approximate the robot's dynamics. 

Among the first model-based approaches to legged locomotion used the Simple Linear Inverted Pendulum (SLIP) model, first proposed by Cavagna \cite{cavagnaMechanicalWorkTerrestrial1977} for modeling locomotion of animals. The SLIP model can be visualized as a point mass on a spring that is passive and conservative, representing the body and leg respectively. With slight modifications, the SLIP model has been used to design impressive locomotion controllers robust to disturbances \cite{raibertTrottingPacingBounding1990, RunningFourLegs, poulakakisStabilityPassiveDynamics2006}. However, the SLIP model isn't directly applicable to all morphologies of robots, including sprawling legged robots such as T-Hex. Furthermore, such approaches usually rely on instantaneous analytical solutions which requires making stric t assumptions regarding the motion of the robot, such as assuming a contact schedule for legs, which can be violated in unstructured terrain.

Modern approaches move away from analytical solutions to optimization-based controllers and use more complex and accurate dynamic models \cite{bledtMITCheetah32018, dicarloDynamicLocomotionMIT2018, kolvenbachTraversingSteepGranular2021, huConstrainedModelPredictive2019, oshinDifferentiableRobustModel2024, grandiaDesignControlBipedal2024, parkOnlinePlanningAutonomous2015, bledtPolicyRegularizedModel, dingRealtimeModelPredictive2019, nosratiRegularizedModelPredictive2025, wieberTrajectoryFreeLinear2006, herdtWalkingThinkingIt2010, prattVirtualModelControl1995, hutterSTARLETHCOMPLIANTQUADRUPEDAL2012, xieIntuitiveApproachQuadruped2015, kocoHybridComplianceControl2019, liuResearchPostureControl2020, torresVirtualModelControl, chenVirtualModelControl2020}. Such controllers plan a center of mass trajectory, typically based on user commands of linear and angular velocity of the body, and then optimize for ground reaction forces exerted by the stance legs to enforce the desired dynamics on the robot's center of mass, such as zero roll and pitch velocity to remain upright or zero vertical acceleration to maintain a base height. Swing legs are typically commanded to follow a reference swing trajectory as is typical with other types of locomotion controllers. 

However, the robot also has to consider how to adjust itself in order to minimize collisions with the walls and ceiling in confined environments. For that, it requires a way of sensing the environment that mere proprioceptive sensors cannot provide. Exteroceptive sensors, such as RGB-D cameras and LiDARs, that provide point clouds containing depth information of the surrounding terrain prove to be a necessity. To understand the environment's geometry, one proposed method is to split the pointcloud into patches and fit local planes using covariance eigenanalysis, yielding slope and roughness \cite{3DSLAMUsing}. This can be expanded upon to include per-cell uncertainty and height leading to a probabilistic elevation map \cite{fankhauserRobotcentricElevationMapping2014, PDFProbabilisticTerrain}. However, these elevation maps do not consider ceilings or overhangs, which introduce multiple z-values. To solve this, the OctoMap \cite{hornungOctoMapEfficientProbabilistic2013} provides full 3D occupancy for the entire space instead of estimating surfaces. On the other hand, Voxblox \cite{oleynikovaVoxbloxIncremental3D2017} utilizes signed distance fields (SDFs) to allow clearance or available space to be obtained efficiently. From there, traversibility; that is, a metric determining safer spots for the robot to go to; is computed. This could range from determining rules for desired slope, roughness, and clearance to learning foothold quality via a network and robot experience. Now, given traversability, the robot must choose how to adapt its configuration or posture to fit in; this can be achieved via abstracting the robot into a deformable bounding box and creating a trajectory of desired body postures \cite{buchananWalkingPostureAdaptation2019, buchananPerceptiveWholebodyPlanning2021}. Perception-based locomotion has also seen significant integration in RL as demonstrated by Miki et al. who used elevation maps as inputs into RL policy \cite{mikiLearningRobustPerceptive2022}. These were replaced with full 3D volumetric representation for confined spaces \cite{mikiLearningWalkConfined2024}. 
